\documentclass[10pt,a4paper]{amsart}
\usepackage[margin=19mm]{geometry}
\usepackage{amsmath,amssymb,mathtools}
\usepackage[scr=rsfs,cal=euler]{mathalfa}
\usepackage{xfrac}
\usepackage[unicode,hidelinks,bookmarks=false]{hyperref}
\newcommand{\ie}{i.e.\ }
\newcommand{\cf}{cf.\ }
\newcommand{\resp}{respectively\ }
\newcommand{\Subsection}{\subsection}
\providecommand{\nobreakdash}{\nobreak\hbox{-}}
\setlength{\emergencystretch}{2em}
\multlinegap0pt
\usepackage{cleveref}
\newcommand{\eps}{\varepsilon}
\newcommand{\redH}{\widetilde{\mathrm H}}
\theoremstyle{plain}
\newtheorem{theorem}{Theorem}[section]
\newtheorem{lemma}[theorem]{Lemma}
\newtheorem{proposition}[theorem]{Proposition}
\newtheorem{corollary}[theorem]{Corollary}
\newtheorem{conjecture}[theorem]{Conjecture}
\theoremstyle{definition}
\newtheorem{definition}[theorem]{Definition}
\newtheorem{example}[theorem]{Example}
\newtheorem{remark}[theorem]{Remark}
\title[Homological Trimming and Regularity of Filtrations via Local Obstruction Modules]{Homological Trimming and Regularity of Filtrations via Local Obstruction Modules}
\author{Siddharth Pritam}
\date{Source: arXiv:2609.28160v1 (CC BY 4.0)}
\begin{document}
\maketitle

\noindent\emph{Source and licence.} The delivered work is the article shown in the title above, version arXiv:2609.28160v1 (CC BY 4.0), distributed under the Creative Commons Attribution 4.0 International licence, as recorded in the arXiv licence metadata for this exact version and shown on the abstract page saved with this item. The full licence notice, which carries the licence URL and the address of that abstract page, is the bundled file \texttt{licenses/arxiv-2609.28160-CC-BY-4.0.txt}.
\par\smallskip

\noindent\textit{Editorial scope.} Counted unit: the article's Proposition (source label prop:no-go) (source0.tex lines 1838--1850, source label \texttt{prop:no-go}): ``Exact slack is not Lipschitz Already for edge-filtered flag complexes, there is no constant for which f- g , r f-r g , L f-L g , R f-R g C f-g holds for every pair of edge functions and every generato''. It is delivered together with the article's complete original proof (source0.tex lines 2294--2339, one proof environment printed later in the article, detached from the statement by the article's own arrangement, 46 lines), copied line for line from the article's own text. Required context retained uncounted: none beyond the counted statement and proof. Editorial formatting only: an AMS \texttt{amsart} preamble that restates the article's own macro definitions and its own theorem declarations, and the theorem environments the retained lines use; the packages cleveref are loaded to match the article's own setup; the article's own class file is neither shipped nor input; the retained environments are renumbered sequentially in order of appearance, whereas the article prints them with its own numbering; the article's equation numbering is not reproduced and every cross-reference inside the retained text resolves to the wrapper's numbering. No citation occurs inside any retained line, so the wrapper carries no bibliography; All retained bodies are exact copies of the bundled \texttt{sources/211556/source0.tex}, so no omission receipt applies. No asset-inclusion command occurs inside any retained span and the package ships no image asset.


\section{The counted unit: Proposition (source label prop:no-go) and its complete original proof}

\begin{proposition}[Exact slack is not Lipschitz]
\label{prop:no-go}
Already for edge-filtered flag complexes, there is no constant $C$ for which
$$
\max\{|\ell_\gamma^f-\ell_\gamma^g|,
       |r_\gamma^f-r_\gamma^g|,
       |L_\gamma^f-L_\gamma^g|,
       |R_\gamma^f-R_\gamma^g|\}
\leq C\|f-g\|_\infty
$$
holds for every pair of edge functions $f,g$ and every generator $\gamma$
that is regular in both filtrations.
\end{proposition}

\begin{proof}[Proof of \Cref{prop:no-go}]
We give an explicit edge-filtration example. Fix numbers $1<s<T$ and let
$0<\eps<s-1$. Take an edge $e=uv$ with value $0$. Its common neighbours are
$$
r,a,b,c,d,h,j.
$$
Give all edges from $u$ or $v$ to these vertices value $-2$. Among the common
neighbours, include the six edges
$$
ra,\ ab,\ bc,\ rd,\ dh,\ hj
$$
with value $-1$. They form two paths sharing the vertex $r$. Add the edges
$cr$, $rb$, and $jr$, and include no other edges among the common neighbours.

Define two edge functions $f$ and $g$ that agree except on $cr$ and $rb$:
$$
f(cr)=f(rb)=s,
\qquad
g(cr)=s-\eps,
\qquad
g(rb)=s+\eps,
$$
and set $f(jr)=g(jr)=T$. Thus $\|f-g\|_\infty=\eps$.

At time $-1$, the filtered link of $e$ is a tree and is therefore acyclic. In
the $f$-filtration, the edge $cr$ closes the first four-cycle at the same time
that the diagonal $rb$ appears. The flag complex immediately contains the two
triangles that fill this cycle, so no obstruction is created. At time $T$, the
edge $jr$ closes the second four-cycle, which has no diagonal and gives a
persistent class in $\redH_1$.

In the $g$-filtration, the first cycle instead exists on
$[s-\eps,s+\eps)$: $cr$ creates it and $rb$ later fills it. The second cycle
still appears at $T$. Consequently, the regularity interval of $e$ containing
its anchor is
$$
[-1,T)
\quad\text{for }f,
\qquad
[-1,s-\eps)
\quad\text{for }g.
$$
The right endpoint and forward slack therefore change by
$T-s+\eps$, while the perturbation size is $\eps$. Letting $\eps$ tend to
zero rules out a universal Lipschitz constant.
\end{proof}

\end{document}
